
Pour vérifier que la suite est exacte, on peut le faire localement sur \(T\), et utiliser les sections définies dans l'énoncé de 2.5. On a alors une base de \(\underline{N}_{T/X}\) formée des images des \(f_i,g_j,h_k\), une base de \(i^*(\underline{N}_{Y/X})\) formée des images des \(f_i,g_j\), et une base de \(j^*(\underline{N}_{Z/X})\) formée des images des \(f_i,h_k\). Il est alors clair que \(i^*(\underline{N}_{Y/X})\times j^*(\underline{N}_{Z/X})\to \underline{N}_{T/X}\) est surjectif, et que son noyau est formé des éléments ayant des coordonnées nulles sur les \((\overline g_j,0)\) et les \((0,\overline h_k)\), et des coordonnées égales pour tout \(i\) sur \((\overline f_i,0)\) et \((0,\overline f_i)\). Par suite, la suite est également exacte au milieu. Or les trois \(\underline{O}_T\)-Modules intervenant sont localement libres de type fini; comme \(\underline{\Tor}_1^{\underline{O}_X}(\underline{O}_Y,\underline{O}_Z)\) est de rang \(p\) d'après la démonstration du i), on voit immédiatement que la suite est exacte.

\subsection*{2.6.}
On suppose maintenant que \(X\) est un schéma ayant un faisceau ample. Tout sous-schéma fermé \(Y\) de \(X\) a alors un faisceau ample. On sait que dans ce cas ``il existe des résolutions globales'', et que les \(K^{\bullet}\) formés respectivement avec les faisceaux localement libres de type fini et les complexes parfaits sont identiques (II, 2.2.9). Si en outre l'immersion \(i:Y\to X\) est régulière, c'est un morphisme parfait, et on peut définir un homomorphisme additif \(i_*:K^{\bullet}(Y)\longrightarrow K^{\bullet}(X)\) (IV 2.12), en associant à la classe d'un complexe parfait \(\underline{L}^\bullet\) la classe du complexe \(\mathbb Ri_*(\underline{L}^\bullet)=i_*(\underline{L}^\bullet)\) (\(i_*\) étant un foncteur exact).

\textbf{Proposition 2.7.} \uline{Soient} \(X\) \uline{un schéma ayant un faisceau ample}, \(i:Y\to X\) \uline{une immersion fermée régulière}, \(\underline{N}\) \uline{le faisceau conormal de} \(Y\) \uline{dans} \(X\). \uline{Alors}, \uline{pour tout} \(x\in K^{\bullet}(Y)\), \uline{on a la relation}:
\begin{equation}\notag
i^*(i_*(x))=x.\lambda_{-1}(N),
\end{equation}
\uline{où} \(N\) \uline{est la classe de} \(\underline{N}\) \uline{dans} \(K^{\bullet}(Y)\), \uline{et} \(\lambda_{-1}(N)=\sum_{i\ge0}(-1)^i\lambda^i(N)\).


